\subsubsection{试验函数}\label{subsec:test-functions}

设 $\Omega$ 为 $\mathbb{R}^m$ 中非空开集. 作如下定义.

\begin{Definition}
\label{def:test-function}
若定义在 $\Omega$ 上的函数 $f$ 满足 $f\in C^\infty(\Omega)$, 且存在紧集 $K\subset\Omega$, 使 $f$ 的支集位于 $K$ 中, 则称其为\textbf{试验函数}. $\Omega$ 上所有试验函数构成的集合记为 $\mathcal{D}(\Omega)=C_0^\infty(\Omega)$.
\end{Definition}

显然, $\mathcal{D}(\Omega)$ 是线性空间. 为进行分析, 需要收敛的概念. 可以在 $\mathcal{D}(\Omega)$ 中定义开集, 并使用一般拓扑的概念. 然而, 对偏微分方程中的大多数目的, 这并非必要; 只需定义序列的收敛. 定义如下.

\begin{Definition}
\label{def:test-function-convergence}
设 $\phi_n$, $n\in\mathbb{N}$ 及 $\phi$ 为 $\mathcal{D}(\Omega)$ 的元素. 若存在 $\Omega$ 的紧子集 $K$, 使所有 $\phi_n$ (以及 $\phi$) 的支集均位于 $K$ 内, 且 $\phi_n$ 以及 $\phi_n$ 的任意阶导数分别一致收敛到 $\phi$ 及其相应导数, 则称 $\phi_n$ 在 $\mathcal{D}(\Omega)$ 中\textbf{收敛}到 $\phi$.
\end{Definition}

\begin{Remark}
\label{rem:test-function-convergence-nonmetric}
注意, 上面定义的收敛概念并不来自度量或范数.
\end{Remark}

知道具有某些性质的试验函数存在通常很重要; 例如, 常常需要一个在给定点 $\bm{y}$ 的小邻域内为正, 在该邻域外为零的函数. 这样的函数可以显式给出:
\begin{equation}\label{eq:smooth-compact-bump-function}
\phi_{\bm{y},\epsilon}(\bm{x})=\begin{cases}\exp\left(-\dfrac{\epsilon^2}{\epsilon^2-|\bm{x}-\bm{y}|^2}\right),&|\bm{x}-\bm{y}|<\epsilon,\\0,&\text{otherwise.}\end{cases}
\end{equation}


事实上, 可以利用这一例子生成其他试验函数的例子. 以下定理指出, 任意具有紧支集的连续函数均可由试验函数一致逼近.

\begin{Theorem}
\label{thm:test-functions-uniform-approximation}
设 $K$ 为 $\Omega$ 的紧子集, 且 $f\in C(\Omega)$ 的支集包含于 $K$. 对 $\epsilon>0$, 令
\begin{equation}\label{eq:test-function-uniform-approximation}
f_\epsilon(\bm{x})=\dfrac1{C(\epsilon)}\int_K\phi_{\bm{y},\epsilon}(\bm{x})f(\bm{y})\,\mathrm{d}\bm{y},
\end{equation}
其中
\begin{equation}\label{eq:test-function-bump-normalization}
C(\epsilon)=\int_{\mathbb{R}^m}\phi_{\bm{y},\epsilon}(\bm{x})\,\mathrm{d}\bm{y}.
\end{equation}
若 $\epsilon<\operatorname{dist}(K,\partial\Omega)$, 则 $f_\epsilon\in\mathcal{D}(\Omega)$; 此外, 当 $\epsilon\to0$ 时, $f_\epsilon\to f$ 一致成立.
\end{Theorem}

证明留作习题.

以类似方式, 可以构造在给定集合上等于 1, 在另一集合上等于 0 的试验函数.

\begin{Theorem}
\label{thm:smooth-cutoff}
设 $K$ 为 $\Omega$ 的紧子集, 且 $U\subset\Omega$ 为包含 $K$ 的开集. 那么存在一个试验函数, 它在 $K$ 上等于 1, 在 $U$ 外等于 0, 并在 $U\setminus K$ 上取值于 $[0,1]$.
\end{Theorem}

\begin{Proof}
取 $\epsilon>0$, 使 $K$ 的 $\epsilon$ 邻域包含于 $U$. 设 $K_1$ 为 $K$ 的 $\epsilon/3$ 邻域的闭包, 并定义
\begin{equation}\label{eq:continuous-neighborhood-cutoff}
f(\bm{x})=1-\min\left\{1,\dfrac3\epsilon\operatorname{dist}(\bm{x},K_1)\right\}.
\end{equation}
函数 $f$ 连续, 在 $K_1$ 上等于 1, 在 $K$ 的 $2\epsilon/3$ 邻域外等于零. 由 \eqref{eq:test-function-uniform-approximation} 定义的 $f_{\epsilon/3}$ 具有定理所要求的性质.
\end{Proof}

偏微分方程中的许多证明涉及归结为在某一点的小邻域内作局部考察. (例如见 Chapter~\ref{chap:linear-elliptic-equations}.) 实现这一归结的工具称为单位分解.

\begin{Definition}
\label{def:partition-of-unity}
设 $U_i$, $i\in\mathbb{N}$ 为 $\Omega$ 的一族有界开子集, 满足
\begin{enumerate}
\item 每个 $U_i$ 的闭包包含于 $\Omega$,
\item $\Omega$ 的每个紧子集只与有限个 $U_i$ 相交 (这一性质称为局部有限性), 并且
\item $\bigcup_{i\in\mathbb{N}}U_i=\Omega$.
\end{enumerate}
从属于覆盖 $\{U_i\}$ 的\textbf{单位分解}是一组试验函数 $\phi_i$, 满足
\begin{enumerate}
\item $0\leq\phi_i\leq1$,
\item $\operatorname{supp}\phi_i\subset U_i$,
\item 对每个 $\bm{x}\in\Omega$, 有 $\sum_{i\in\mathbb{N}}\phi_i(\bm{x})=1$.
\end{enumerate}
\end{Definition}

以下定理指出, 这样的单位分解存在.

\begin{Theorem}
\label{thm:partition-of-unity-existence}
设集合族 $U_i$, $i\in\mathbb{N}$ 具有 Definition~\ref{def:partition-of-unity} 中陈述的性质. 那么存在从属于覆盖 $\{U_i\}$ 的单位分解.
\end{Theorem}

\begin{Proof}
首先构造新的覆盖 $\{V_i\}$, 其中 $V_i$ 具有 Definition~\ref{def:partition-of-unity} 的全部性质, 并且 $V_i$ 的闭包包含于 $U_i$. 通过归纳构造 $V_i$. 假设已经找到 $V_1$, $V_2$, $\ldots$, $V_{k-1}$, 使 $U_j$ 包含 $\overline V_j$, 且
\begin{equation}\label{eq:shrinking-cover-inductive-invariant}
\Omega=\bigcup_{j=1}^{k-1}V_j\cup\bigcup_{j=k}^\infty U_j.
\end{equation}
设 $F_k$ 为集合
\begin{equation}\label{eq:shrinking-cover-complement-set}
\bigcup_{j=1}^{k-1}V_j\cup\bigcup_{j=k+1}^\infty U_j.
\end{equation}
的补集. 那么 $F_k$ 为包含于 $U_k$ 的闭集. 选取任意包含 $F_k$ 且满足 $\overline V_k\subset U_k$ 的开集作为 $V_k$. 每个点 $\bm{x}\in\Omega$ 只包含于有限个 $U_i$; 因而存在 $N\in\mathbb{N}$, 使 $\bm{x}\notin\bigcup_{j=N+1}^\infty U_j$. 但这意味着 $\bm{x}\in\bigcup_{j=1}^NV_j$. 因此 $V_i$ 具有 Definition~\ref{def:partition-of-unity} 的性质 3; 另外两个性质由 $V_i\subset U_i$ 显然得出.

设 $W_k$ 为开集, 满足 $\overline V_k\subset W_k$, $\overline W_k\subset U_k$. 根据 Theorem~\ref{thm:smooth-cutoff}, 存在试验函数 $\psi_k$, 在 $\overline V_k$ 上等于 1, 在 $W_k$ 外等于零, 其余处的值介于 0 与 1 之间. 令
\begin{equation}\label{eq:partition-of-unity-normalizing-sum}
\psi(\bm{x})=\sum_{k\in\mathbb{N}}\psi_k(\bm{x}).
\end{equation}
由于 Definition~\ref{def:partition-of-unity} 的性质 2, 在任意给定 $\bm{x}$ 的邻域内, \eqref{eq:partition-of-unity-normalizing-sum} 的右端只有有限个非零项, 因而不存在收敛问题. 函数 $\phi_k:=\psi_k/\psi$ 给出所需的单位分解.
\end{Proof}
